%  =============================================================================
%  Datei-Hinweise & Konfiguration: siehe config.md
%  -> Abschnitt "anhang/ki-prompts/tool-claude.tex"
%  =============================================================================

\section{Anthropic Claude (Opus~5)}
\label{sec:prompts-claude}

\kiwerkzeug{Anthropic Claude, Modell Opus 5}{Modellbezeichner
\texttt{claude-opus-5}, veröffentlicht am 24.07.2026}{https://claude.ai}
% Eintrag im KI-Verzeichnis (Leitfaden Kap. 6.2.4) erzwingen mittels \nocite.
\nocite{claudeOpus2026}

Eingesetzt für die Verzeichnisse des Vorspanns und den Aufbau des Anhangs
sowie für zwei Werkzeuge der Projektumgebung: die Anpassung des
Extraktionsskripts, das Quellcode- und Datenbestände der Google-Cloud-Umgebung
für das elektronische Archiv aufbereitet, und die Erzeugung der Testskripte
zur Ergebniskontrolle.

\subsection{Erstellung des Glossars}
\label{sec:prompts-claude-glossar}
\begin{enumerate}
    \item \enquote{Formuliere zu den folgenden Fachbegriffen jeweils eine
            prägnante, allgemeinverständliche Glossardefinition von höchstens
            zwei Sätzen im sachlichen wissenschaftlichen Stil: [\dots]}
            \par\hfill\textit{Kontext: Glossar (Vorspann)}

    \item \enquote{Prüfe die folgenden Glossareinträge auf fachliche
            Korrektheit, Widerspruchsfreiheit und einheitliche Formulierung:
            [\dots]}
            \par\hfill\textit{Kontext: Glossar (Vorspann)}
\end{enumerate}

\subsection{Erstellung des Formelzeichenverzeichnisses}
\label{sec:prompts-claude-formelzeichen}
\begin{enumerate}
    \item \enquote{Extrahiere aus den folgenden nummerierten Gleichungen
            sämtliche Formelzeichen und formuliere je Zeichen eine
            Kurzerklärung, die Bedeutung und Wertebereich beziehungsweise
            Einheit benennt: [\dots]}
            \par\hfill\textit{Kontext: Formelzeichenverzeichnis (Vorspann)}

    \item \enquote{Prüfe, ob jedes in den Gleichungen verwendete Formelzeichen
            im Verzeichnis geführt wird und ob umgekehrt jeder Eintrag des
            Verzeichnisses im Text tatsächlich vorkommt; nenne Abweichungen in
            beide Richtungen: [\dots]}
            \par\hfill\textit{Kontext: Formelzeichenverzeichnis (Vorspann)}
\end{enumerate}

\subsection{Aufbau des Anhangs}
\label{sec:prompts-claude-anhang}
\begin{enumerate}
    \item \enquote{Schlage eine Gliederung für einen dreiteiligen Anhang vor,
            der ein elektronisches Quellcodearchiv, die Belege einer
            Evaluation und die verwendeten KI-Prompts umfasst; benenne je
            Teil, welche Angaben erforderlich sind, damit ein Dritter die
            Ergebnisse nachvollziehen kann: [\dots]}
            \par\hfill\textit{Kontext: Anhang~\ref{cha:anhang-archiv} bis
            Anhang~\ref{cha:anhang-prompts}}

    \item \enquote{Formuliere zu den nachfolgenden Rohprotokollen jeweils
            einen einleitenden Absatz, der Herkunft, Erzeugungszeitpunkt und
            Aussagegehalt des Belegs benennt, ohne die Zahlen selbst zu
            wiederholen: [\dots]}
            \par\hfill\textit{Kontext:
            Anhang~\ref{sec:anhang-evaluationsbelege}}
\end{enumerate}

\subsection{Anpassung des Extraktionsskripts für Code- und Datenquellen}
\label{sec:prompts-claude-extraktion}

Das Skript zur Aufbereitung des Quellcodearchivs lag aus vorangegangenen
Projekten bereits vor; Gegenstand der Prompts war ausschließlich seine
Anpassung an die Zielarchitektur dieser Arbeit (\cref{fig:zielarchitektur})
mit ihren zwei virtuellen Maschinen, Journaldateien und zeitgesteuerten
Aufträgen.

\begin{enumerate}
    \item \enquote{Erweitere das vorhandene Extraktionsskript so, dass es
            neben Python- und Shell-Quelltexten auch die
            \texttt{systemd}-Einheiten und Cron-Definitionen beider
            virtueller Maschinen erfasst und je Datei Herkunftsmaschine sowie
            Ausführungsintervall in einem Manifest festhält: [\dots]}
            \par\hfill\textit{Kontext:
            Anhang~\ref{sec:anhang-archiv-quellcode}}

    \item \enquote{Ergänze das Skript um eine Redaktionsstufe, die
            Zugangsdaten sowie Konto- und Ordernummern aus den Journaldateien
            entfernt, bevor diese in das Archiv aufgenommen werden, und die
            jede Ersetzung protokolliert: [\dots]}
            \par\hfill\textit{Kontext: Anhang~\ref{sec:anhang-archiv-daten}}

    \item \enquote{Stelle sicher, dass die erzeugte Archivdatei bei
            unverändertem Inhalt reproduzierbar dieselbe Prüfsumme liefert,
            und erläutere, welche Zeitstempel dafür festzusetzen sind:
            [\dots]}
            \par\hfill\textit{Kontext:
            Anhang~\ref{sec:anhang-archiv-identitaet}}
\end{enumerate}

\subsection{Erzeugung der Testskripte zur Ergebniskontrolle}
\label{sec:prompts-claude-tests}
\begin{enumerate}
    \item \enquote{Entwirf zu der nachfolgenden Auswertungsfunktion eine
            Testsuite mit gemockten Schnittstellen, die jede
            sicherheitsrelevante Verzweigung abdeckt -- insbesondere die
            Pfade, die zu einer Ablehnung oder zu einem Handelsstopp führen:
            [\dots]}
            \par\hfill\textit{Kontext: \cref{sec:umsetzung-implementierung}}

    \item \enquote{Formuliere Testfälle, die prüfen, ob die im Protokoll
            ausgewiesenen Kennzahlen aus denselben Eingangsdaten
            reproduzierbar hervorgehen, und benenne die Toleranz, innerhalb
            derer Gleitkommaabweichungen zulässig sind: [\dots]}
            \par\hfill\textit{Kontext: \cref{sec:umsetzung-evaluation}}
\end{enumerate}

%  =============================================================================
%  Ende
%  =============================================================================
